%TO CHANGE THE TITLE OR THE FRONT PAGE IN GENERAL GO TO LINE 65


\documentclass{article}
\usepackage{header}
%\usepackage{codespace}
%\usepackage{booktabs}


% Sub-preambles
% https://github.com/MartinScharrer/standalone

% Encodings
\usepackage{amsmath,amssymb,gensymb,textcomp,unicode-math}

% Better tables
% Wide tables go to https://tex.stackexchange.com/q/332902
\usepackage{array,multicol,multirow,siunitx,tabularx}

% Better enum
\usepackage{enumitem}
\usepackage{slashbox}

% Graphics
\usepackage{caption}
\usepackage{sidecap}
\usepackage{float}
\usepackage{geometry}
%\usepackage{tikz} %per creare simboli
\usepackage{afterpage} % Per eseguire il codice solo dopo la prima pagina
\usepackage{changepage}


\usepackage{xcolor}    % Per definire il colore desiderato
\definecolor{pantone}     {HTML}{9B0014}
\definecolor{rosso-scuro} {HTML}{800020}
\definecolor{rosso-nasa}  {HTML}{FC3D21}
\definecolor{blu-scuro}   {HTML}{001F3F}
\definecolor{blu-chiaro}  {HTML}{1F4E79}
\definecolor{verde}      {HTML}{006400}
\definecolor{marrone}     {HTML}{654321}
\definecolor{oro}         {HTML}{FFD700}
\definecolor{bianco}      {HTML}{FFFFFF}
\definecolor{nero}        {HTML}{000000}

\captionsetup{font=small}
\usepackage{fontspec}
\setmainfont{Bauer-Bodoni-Regular.otf}

% \usepackage[inkscapeformat=png]{svg}
% Allow setting >max< width of figure
% 'export' allows adjustbox keys in \includegraphics
\usepackage[export]{adjustbox}




% For demonstration purposes, remove in production
\usepackage{mwe}

% Configurations
\newcounter{memberrowno}
\setcounter{memberrowno}{0}

%%%%%%%%%%%%%%%%%%%%%%%%%%%%%%%%%%%%%%%%%%%%%%%%%%
\ocoursename{DIPARTIMENTO DI FISICA E ASTRONOMIA}
\oreportcity{\hspace{2.5cm} PADOVA}%\hspace{2.5cm}  its the white space that must be left to align the city yer and group number
\ocoursenamedue{GALILEO GALILEI}
\oreporttype{SPERIMENTAZIONI DI FISICA 1}
\oreportyear{\hspace{7.05cm} A.A. 22-23}
\ogroupnumber{\hspace{4.565cm} GXX}

\otitlelineuno{WRITE}
\otitlelinedue{REPORT}
\otitlelinetre{NAME}
\otitlelinefour{HERE}
%%%%%%%%%%%%%%%%%%%%%%%%%%%%%%%%%%%%%%%%%%%55





\title{}


\reportlayout%

% Custom commands
\newcommand*\mean[1]{\bar{#1}}


\begin{document}
% margini per la prima pagina
\newgeometry{left=0cm,right=0cm,top=0cm,bottom=0cm}
\coverpage%
%imposta i bordi per le pagine suggessive
\newgeometry{
    paper=a4paper,
    vmargin=2cm,
    hmargin=2cm,
    %top=0cm,
    includeheadfoot=true,
    headheight=30pt
  }
% Esegui il codice solo sulla prima pagina
% Imposta lo sfondo bianco per tutte le pagine

\color{black}
  \thispagestyle{empty}
  %\setmainfont{Roboto}
  ~
  \vspace{17cm}
  \begin{figure}[H]
      \begin{minipage}{0.55\textwidth}
    This work is licensed under a \href{https://creativecommons.org/licenses/by-nc-sa/3.0/deed.en}{Creative Commons "Attribution-NonCommercial-ShareAlike 3.0 Unported"} license.
    \end{minipage}%
    \begin{minipage}{0.4\textwidth}
    \centering
    \href{https://creativecommons.org/licenses/by-nc-sa/3.0/deed.en}{
    \includegraphics[width=.7\textwidth]{graphics/header/ccc.png}
    }
    \end{minipage}
  \end{figure}
   
\clearpage




\tableofcontents
%ingrandisci la tabella sotto e riformattala
\vfill
\begin{center}
    \begin{tabular}{lr}ㅤㅤ
    \color{White} aaaaaaaaaaaaaaaaaaaaa& \textbf{Gruppo XX}\\
     \end{tabular}
    \newline
    \begin{tabular}{r|l|l}
      nome cognome1 & nome.cognome1@studenti.unipd.it & xxxxxxx   \\\hline
      nome cognome2 & nome.cognome2@studenti.unipd.it & xxxxxxx\\\hline
      nome cognome3 & nome.cognome3@sutdenti.unipd.it & xxxxxxx   \\
      % Student: & Name      & ID \\
    \end{tabular}
\end{center}


\newpage



\input{chapters/cap1.tex}
\input{chapters/cap2.tex}
\input{chapters/cap3.tex}
\input{chapters/cap4.tex}

\input{chapters/appendice.tex}


\nocite{*}

\end{document}
